\documentclass{article}

\usepackage{polski}
\usepackage[utf8]{inputenc}
\usepackage{graphicx}
\usepackage{xcolor}
\usepackage{float}
\usepackage{caption}
\usepackage{array}
\usepackage{pbox}

\newcommand\tab[1][1cm]{\hspace*{#1}}

\title{Dokumentacja projektowa}
\date{2018-03-18}
\author{Jędrzej Kozal}

\begin{document}

\begin{titlepage}
	\centering
	\includegraphics[width=0.25\textwidth]{logo_pol_wroclaw.png}\par\vspace{1cm}
	{\scshape\LARGE Politechnika Wrocławska \par}
	\vspace{1cm}
	{\scshape\Large Podstawy Obliczeń Neuronowych\par}
	\vspace{1.5cm}
	{\huge\bfseries Sieci radialne w predykcji procesów medycznych - choroba Parkinsona \par}
	\vspace{2cm}
	{\Large\itshape Filip Guzy\par}
	{\Large\itshape Jędrzej Kozal\par}

	\vfill
	prowadzący\par
	Dr inż.~Edward \textsc{Puchała}

	\vfill

% Bottom of the page
	{\large \today\par}
\end{titlepage}

\tableofcontents
\newpage


\section{Wstęp}

\subsection{Cel projektu}

Celem projektu było zbadanie skuteczności radialnych sieci neuronowych do predykcji przebiegu choroby Parkinsona.

\subsection{Wykorzystane narzędzia}

Do realizacji projektu wykorzystano środowisko Matlab, oraz dane z zbioru uczącego UCI Oxford Parkinson's Disease Telemonitoring Dataset, udostępnione przez autorów \cite{zbior_uczacy}.

\subsection{Aspekt medyczny}

Coś o Parkinsonie i danych.

\section{Zarys podstaw teoretycznych}

\section{Osiągnięte rezultaty}

\section{Wnioski}

\newpage
\begin{thebibliography}{9}

\bibitem{zbior_uczacy}
A Tsanas, MA Little, PE McSharry, LO Ramig (2009) 
'Accurate telemonitoring of Parkinson’s disease progression by non-invasive speech tests', 
IEEE Transactions on Biomedical Engineering (to appear). 
\\\texttt{https://archive.ics.uci.edu/ml/datasets/Parkinsons+Telemonitoring}

\bibitem{zbior_uczacy_1}
Little MA, McSharry PE, Hunter EJ, Ramig LO (2009), 
'Suitability of dysphonia measurements for telemonitoring of Parkinson's disease', 
IEEE Transactions on Biomedical Engineering, 56(4):1015-1022 

\bibitem{zbior_uczacy_2}
Little MA, McSharry PE, Roberts SJ, Costello DAE, Moroz IM. 
'Exploiting Nonlinear Recurrence and Fractal Scaling Properties for Voice Disorder Detection', 
BioMedical Engineering OnLine 2007, 6:23 (26 June 2007) 

\bibitem{Osowski}
Stanisław Osowski.
\textit{Sieci Neuronowe w ujęciu algorytmicznym}
Wydawnictwo Naukowo-Techniczne, Warszawa 1996

\bibitem{Osowski}
Stanisław Osowski.
\textit{Sieci Neuronowe do przetwarzania informacji}
Oficyna Wydawnicza Politechniki Warszawskiej, Warszawa 2006

\bibitem{matlab}
Oficjalna dokumentacja Mathworks Deep Learning Toolbox. 
\\\texttt{https://www.mathworks.com/help/deeplearning/examples/radial-basis\\-approximation.html}

\bibitem{owards Data Sience}
Towards Data Science blog
\\\texttt{https://towardsdatascience.com/radial-basis-functions-neural-\\networks-all-we-need-to-know-9a88cc053448}


\end{thebibliography}
\newpage

\listoffigures
\newpage

\end{document}